% Part: second-order-logic
% Chapter: metatheory
% Section: introduction

\documentclass[../../../include/open-logic-section]{subfiles}

\begin{document}

\olfileid{sol}{met}{int}

\olsection{ভূমিকা}

প্রথম-ক্রমের যুক্তিবিদ্যার কয়েকটি সুন্দর ধর্ম আছে। আমরা জানি, সেটি
নির্ণেয় নয়, তবে অন্তত স্বতঃসিদ্ধায়নযোগ্য। অর্থাৎ, প্রথম-ক্রমের
যুক্তিবিদ্যার জন্য এমন বিশুদ্ধ ও পূর্ণ প্রমাণ-ব্যবস্থা আছে, যা এমন একটি
!!a{derivability} সম্পর্ক~$\Proves$ দেয় যে, !!{sentence}s-এর যেকোনো
সেট~$\Gamma$ ও !!{sentence}~$!A$-এর জন্য $\Gamma \Entails !A$ যদি এবং
কেবল যদি $\Gamma \Proves !A$। বিশেষত, এর অর্থ প্রথম-ক্রমের যুক্তিবিদ্যার
বৈধ বাক্যগুলি !!{computably
  enumerable}। এমন একটি গণনসাধ্য অপেক্ষক~$f\colon \Nat \to
\Sent[L]$ আছে, যেখানে~$f$-এর মানগুলি ঠিক~$\Lang{L}$-এর বৈধ !!{sentence}s। এর
কারণ, !!{derivation}s তালিকায়িত করা যায়, এবং যেগুলি একটি মাত্র
!!{sentence}~!!{derive} করে সেগুলিকে সেই !!{sentence}-এ পাঠানো যায়।
দ্বিতীয়-ক্রমের যুক্তিবিদ্যা প্রথম-ক্রমের যুক্তিবিদ্যার চেয়ে বেশি
প্রকাশক্ষম, তাই সাধারণভাবে তার বৈধ বাক্যগুলি ধারণ করা আরও জটিল। বস্তুত,
আমরা দেখাব যে দ্বিতীয়-ক্রমের যুক্তিবিদ্যা শুধু অণির্ণেয়ই নয়, তার বৈধ
বাক্যগুলি এমনকি !!{computably
  enumerable}-ও নয়। এর অর্থ, দ্বিতীয়-ক্রমের যুক্তিবিদ্যার কোনো বিশুদ্ধ
ও পূর্ণ প্রমাণ-ব্যবস্থা থাকতে পারে না (যদিও বিশুদ্ধ কিন্তু অপূর্ণ
প্রমাণ-ব্যবস্থা আছে এবং সেগুলি সত্যিই গবেষণার গুরুত্বপূর্ণ বিষয়)।

প্রথম-ক্রমের যুক্তিবিদ্যার আরও দুটি ধর্ম আছে: সেটি সংহত
(!!{sentence}s-এর কোনো সেট~$\Gamma$-র প্রতিটি সসীম উপসেট
পরিতৃপ্তিযোগ্য হলে $\Gamma$ নিজেও পরিতৃপ্তিযোগ্য), এবং তার ক্ষেত্রে
লোয়েনহাইম--স্কোলেম উপপাদ্য সত্য ($\Gamma$-র একটি অসীম মডেল থাকলে তার
একটি !!a{denumerable} মডেলও আছে)। এই দুটি ফলই দ্বিতীয়-ক্রমের
যুক্তিবিদ্যায় ব্যর্থ। এখানেও কারণটি হলো, দ্বিতীয়-ক্রমের যুক্তিবিদ্যা
!!{domain}s-এর আকার সম্বন্ধে এমন তথ্য প্রকাশ করতে পারে যা প্রথম-ক্রমের
যুক্তিবিদ্যা পারে না।

\end{document}
